\section{Background and Related Work}

\paragraph{BabyLM.}
BabyLM evaluates language models trained under constrained data budgets, with particular emphasis on whether small models can acquire linguistic generalizations from limited input \citep{warstadt2023babylm,hu2024babylm,choshen2026babylm}. The 2026 Strict-Small track limits pretraining data to 10M words or less and asks participants to implement submissions using Hugging Face Transformers. The shared-task setting is useful for this paper because it makes the central constraint explicit: performance should improve by using limited data more effectively, not by silently increasing corpus size or model scale.

The recent BabyLM literature shows that strong results often come from changing several parts of the training recipe at once. The second BabyLM findings report that the best-performing submissions tended to modify the training data, objective, or architecture, and that a hybrid causal-masked architecture was especially successful \citep{hu2024babylm}. GPT-BERT merges causal and masked objectives in a single transformer stack and outperforms masked-only or causal-only variants in the 2024 challenge setting \citep{charpentier2024gptbert}; AntLM similarly alternates causal and masked objectives and attention masks to combine the strengths of CLM and MLM \citep{yu2024antlm}. Other recent BabyLM systems explore orthogonal sources of efficiency, including linear or lightweight token mixers, optimizer changes, curated corpora, masked diffusion objectives, and frequency-informed training \citep{haller2025sample,kosmopoulou2025masked}. These systems are important points of comparison, but they answer a broader question than ours. We ask a narrower question: if the data, tokenizer, architecture, optimizer, schedule, training length, and total MLM mask rate are held fixed, can changing only the masking policy improve sample efficiency?

\paragraph{Why Mask Selection Matters.}
In masked language modeling, the model only receives a learning signal at the positions selected for prediction. The corpus may contain 10M words, but a 15\% mask rate exposes the loss directly on only a fraction of token positions per epoch. Under web-scale pretraining this inefficiency may be tolerable because the model sees enormous redundancy. Under BabyLM Strict-Small constraints, however, the prediction budget is much tighter: repeatedly spending masks on very easy function words, isolated punctuation-adjacent tokens, or one-off examples may be less useful than spending some of the same budget on positions that stress lexical, discourse, or morphological generalization.

\paragraph{Adaptive Masking.}
Masked language modeling usually samples prediction targets uniformly from valid token positions. This is an attractive default because it is simple, unbiased with respect to token position, and does not require task-specific annotation. The downside is that random masking treats all valid positions as equally valuable prediction targets. Prior work has therefore explored masking policies that change what gets predicted. Concept-based curriculum masking evaluates token difficulty with linguistic and knowledge-graph criteria, then masks easy items first and related items later \citep{lee2022efficient}. Scheduled masking changes the masking ratio and masked content over training time, arguing that time-invariant MLM settings are unlikely to be optimal \citep{yang2023learning}. Self-Evolution learning focuses on informative or under-explored tokens and combines adaptive masking with token-specific label smoothing \citep{zhong2023self}. Most directly related to our follow-up experiments, \citet{edman2025mask} adapt MLM masking probabilities in BabyLM according to the model's ability to predict tokens and also use subtoken embeddings to improve morphological generalization.

HHM and Accuracy-Morph are in this adaptive-masking family, but the experimental emphasis is different. We keep the total mask rate fixed at 15\% for every condition and avoid adding a new objective, label smoother, subtoken embedding module, external knowledge graph, or data-selection step. The only intended difference from the random baseline is the collator's target-selection policy. This makes the experiment less likely to produce a large leaderboard gain, but it makes the causal interpretation cleaner: any change in validation loss or official diagnostics is attributable to which positions were predicted, not to a larger training signal.

\paragraph{Why Heavy Hitters.}
The first hypothesis behind HHM is that repeated strings in a local context are often linguistically important. In child-directed or narrative-like text, repeated names, pronouns, and content words help establish discourse participants and event structure. A random mask may hit these positions occasionally, but a heavy-hitter mechanism can deliberately spend part of the fixed mask budget on recurring entity-like mentions. This is meant to encourage the model to learn from local recurrence and to support entity-tracking behavior without adding a parser, a coreference system, or external supervision.

\paragraph{Why Hard Examples.}
The second hypothesis is curriculum-like: tokens that the current model repeatedly predicts poorly may be worth revisiting. The online error component of HHM records high-loss masked targets during training and later makes matching token strings more likely to be masked. This creates a lightweight feedback loop in which the model's own errors influence future prediction targets. The risk is that raw loss can be noisy: rare words, tokenization artifacts, or inherently ambiguous contexts may receive attention without producing useful generalization. This risk motivated the later Accuracy-Morph pilot, which replaces raw loss with smoothed prediction-correctness statistics.

\paragraph{Streaming Sketches.}
Streaming sketches are compact data structures that approximate properties of a stream, such as heavy hitters and frequency moments \citep{cormode2005improved,metwally2005spacesaving}. They are well matched to this setting because the masking policy observes a long stream of token events but should not store unbounded per-token histories. We use Space-Saving sketches to maintain bounded-memory summaries of repeated local entity strings and high-loss target tokens. The follow-up CMS-Morph pilots use Count-Min Sketches to track token and character-trigram difficulty signals. In both cases, the sketch is a small control mechanism over mask selection, not an added source of text or labels.

\paragraph{Morphological Motivation.}
The CMS-Morph and Accuracy-Morph pilots were tested because token-level difficulty may be too narrow in a small-data regime. If the model struggles with one word form, related forms may share prefixes, suffixes, stems, or orthographic patterns. Character trigrams are a deliberately simple proxy for these subtoken regularities: they require no morphological analyzer, work with the same tokenizer, and can generalize a difficulty signal across related surface forms. This connects to the BabyLM finding that subtoken information can improve morphological generalization \citep{edman2025mask}, but our implementation uses character-trigram sketches only to decide future mask positions. It does not change the input representation or add new embeddings.

\paragraph{Design Constraint.}
The central design constraint is that any observed difference should be attributable to mask selection rather than to extra data, a stronger model, a different tokenizer, or a larger prediction budget. This makes the experiment intentionally conservative. A positive result would indicate that a cheap online policy can make the same training budget more useful; a negative or mixed result still identifies which adaptive signals are too noisy or too weak under strict controls. This is also what makes the paper different from many successful BabyLM submissions: rather than proposing a stronger complete system, we isolate one low-cost mechanism and report both where it fails and where a corrected version begins to help.
